\documentclass{article}

\usepackage[nonatbib]{neurips}
\usepackage[numbers]{natbib}

\input{extra_pkgs}

\usepackage{physics}
\usepackage{mathtools}
\usepackage{multirow}
\usepackage{subcaption}
\usepackage{amssymb}
\usepackage{bm}
\usepackage{xspace}

\newcommand{\method}{REIN\xspace}
\newcommand{\eg}{e.g.,\xspace}
\newcommand{\ie}{i.e.,\xspace}
\newcommand{\etal}{\textit{et al.}\xspace}

\title{REIN: Residual Modulation of Time Series Forecasts\\ via LLM-Refined News Regime Signals}

\author{
    Anonymous Author(s)
}

\begin{document}

\maketitle

%% ============================================================
%% ABSTRACT
%% ============================================================
\begin{abstract}
We introduce \method (\textbf{RE}sidual \textbf{I}nformation from \textbf{N}ews), a two-stage framework that augments time series forecasting with LLM-refined textual signals through multiplicative residual modulation.
Given a frozen base forecaster and a corpus of news articles, \method first distills each article into a structured regime vector using a large language model, then aggregates these vectors into a daily regime bank.
A PatchTST-based delta module decomposes the base model's residual into slow (level), shape (intra-day pattern), and spike (extreme event) components, whose magnitudes are modulated---but never additively injected---by the news regime representation.
This multiplicative-only coupling provides a formal guarantee: when no relevant news exists, the modulation knobs collapse to identity and the system gracefully degrades to the base forecaster with bounded perturbation.
We evaluate \method on five diverse datasets spanning Australian electricity load, Australian road traffic, European natural gas demand, Bitcoin hourly prices, and Australian electricity price, paired with both curated domain news and the open GDELT global news corpus.
Each base backbone is chosen per dataset from a pool of PatchTST, DLinear, NLinear, and MLP variants, and evaluated over horizons ranging from 48 to 720 steps.
A counterfactual evaluation protocol with news-blanking and date-permutation tests is used to verify that any improvement is causally attributable to news content rather than to statistical artifacts or spurious correlations with calendar effects.
\end{abstract}

%% ============================================================
%% 1. INTRODUCTION
%% ============================================================
\section{Introduction}
\label{sec:intro}

Accurate time series forecasting underpins decision-making in energy markets, transportation, and finance.
While modern deep forecasters~\citep{PLACEHOLDER_patchtst_nie2023, PLACEHOLDER_dlinear_zeng2023, PLACEHOLDER_informer_zhou2021} achieve strong performance by exploiting temporal patterns, they remain blind to exogenous textual information---weather warnings, policy announcements, supply disruptions---that human analysts routinely incorporate.
Integrating such unstructured text into numerical forecasters is challenging for three reasons:
(1)~raw news is noisy, with most articles irrelevant to the target series;
(2)~the causal link from text to time series varies across domains and time scales;
(3)~naive integration risks \emph{news hallucination}, where the model fabricates spurious correlations that degrade forecasts when news is absent.

Recent work has explored incorporating textual information into time series models through various means: prompt-based LLM forecasters~\citep{PLACEHOLDER_llmtime_gruver2024, PLACEHOLDER_timellm_jin2024}, cross-modal fusion~\citep{PLACEHOLDER_crossformer_zhang2023}, and retrieval-augmented approaches~\citep{PLACEHOLDER_rag_ts_2024}.
However, these approaches typically treat text as an additive input feature, offering no formal control over its influence and no guarantee of graceful degradation.

We propose \method, a framework built on three design principles that address these challenges:

\begin{itemize}[leftmargin=1.5em, itemsep=2pt]
    \item \textbf{Structured news distillation.}
    An LLM refines each raw article into a fixed-schema regime vector (5~continuous dimensions + 15~topic tags + a 384-d text embedding), filtering noise while preserving actionable signals.

    \item \textbf{Multiplicative-only modulation.}
    News signals \emph{never} enter the residual additively.
    Instead, three bounded modulation knobs---$\lambda_{\text{base}}$ (trust), $g_{\text{shape}}$ (shape gain), and $b_{\text{spike}}$ (spike bias)---\emph{scale} the time-series-derived residual components.
    When no news is active, all knobs collapse to identity values ($\lambda \to \lambda_{\min}$, $g \to 1$, $b \to 0$), yielding a formal degradation bound.

    \item \textbf{Counterfactual evaluation protocol.}
    Each test sample is evaluated under three conditions---normal, news-blanked (all-zero regime), and date-permuted (shuffled news-date alignment)---enabling causal attribution of forecast improvements to news content.
\end{itemize}

We evaluate \method on five datasets: Australian electricity load (NSW), Australian road traffic counts (NSW, permanent stations), European natural gas demand (Netherlands), Bitcoin hourly prices, and Australian electricity price (NSW).
Each dataset is paired with a domain-appropriate news corpus---curated energy/renewables outlets for load and price, CBS archival news for gas demand, and the open GDELT global news corpus~\citep{PLACEHOLDER_gdelt_leetaru2013} for traffic and Bitcoin---and is forecast over horizons up to 720 steps.
Across positive-gain datasets, top-decile help rates routinely exceed 75\%, indicating that the delta module is most active precisely where the base forecaster struggles.
The price dataset reveals an instructive failure mode: heavy-tailed distributions amplify rare errors, causing overall MAE to worsen despite a majority of samples being improved---an insight we formalize as the \emph{help-rate/MAE decoupling} phenomenon.


%% ============================================================
%% 2. RELATED WORK
%% ============================================================
\section{Related Work}
\label{sec:related}

\paragraph{Deep time series forecasting.}
Transformer-based models have become dominant in long-horizon forecasting.
PatchTST~\citep{PLACEHOLDER_patchtst_nie2023} introduced channel-independent patching with self-attention, achieving state-of-the-art on standard benchmarks.
DLinear~\citep{PLACEHOLDER_dlinear_zeng2023} demonstrated that simple linear models with trend-seasonal decomposition can match Transformer performance.
iTransformer~\citep{PLACEHOLDER_itransformer_liu2024} inverted the attention dimension to capture variate correlations.
Our base backbone leverages MLP or DLinear architectures, while the delta module employs PatchTST-style encoding---separating ``what the base can learn'' from ``what news can correct.''

\paragraph{Text-augmented forecasting.}
Incorporating textual signals for time series prediction has a long history in finance~\citep{PLACEHOLDER_sentiment_finance_2020} and energy~\citep{PLACEHOLDER_news_energy_2022}.
LLMTime~\citep{PLACEHOLDER_llmtime_gruver2024} repurposed LLMs as zero-shot forecasters by tokenizing numerical values, while Time-LLM~\citep{PLACEHOLDER_timellm_jin2024} aligned time series patches with language model representations.
UniTime~\citep{PLACEHOLDER_unitime_liu2024} unified multiple forecasting domains with language-guided instructions.
These methods either require the LLM at inference time (costly for high-frequency series) or fuse text features additively without controlling their influence.
\method differs by using the LLM \emph{only offline} for news distillation and enforcing strictly multiplicative coupling at inference.

\paragraph{Regime-switching and conditional models.}
Regime-switching models~\citep{PLACEHOLDER_hamilton1989} capture structural breaks by conditioning on latent states.
Recent neural variants learn regime representations from auxiliary data~\citep{PLACEHOLDER_neural_regime_2023}.
Our regime bank is conceptually similar but derives regimes from \emph{explicit} textual signals rather than latent statistical patterns, providing interpretable regime attributions and enabling the counterfactual evaluation protocol.


%% ============================================================
%% 3. METHOD
%% ============================================================
\section{Method}
\label{sec:method}

Figure~\ref{fig:architecture} provides an overview of \method.
Given a univariate time series $\{y_t\}$ and an associated news corpus $\mathcal{N}$, the framework operates in two stages: base training (Stage~1) and delta training (Stage~2).

\begin{figure}[t]
    \centering
    % TODO: Replace with actual architecture figure
    \fbox{\parbox{0.95\textwidth}{\centering\vspace{1.5em}
    \small
    \textbf{Stage 1 (frozen):} \quad History $\bm{h}_z$ $\;\to\;$ Base Backbone $\;\to\;$ $\hat{\bm{y}}^{\text{base}}_z,\; \bm{f}_{\text{base}}$\\[1em]
    \textbf{Offline:} \quad News Corpus $\;\xrightarrow{\text{LLM refine}}\;$ Refined JSONL $\;\xrightarrow{\text{aggregate}}\;$ Regime Bank $\mathcal{B}$\\[1em]
    \textbf{Stage 2:} \quad $\bm{h}_z + \bm{f}_{\text{base}}$ $\;\to\;$ PatchTST Encoder $\;\to\;$ $\bm{T}, \bm{s}$\\
    $\mathcal{B}(d)$ $\;\to\;$ Regime Projector $\;\to\;$ $\bm{v}_{\text{regime}}$\\[0.5em]
    $\bm{s} \to$ \{Slow, Shape, Spike\} Heads $\;\xrightarrow[\text{modulate}]{\lambda,\, g,\, b}\;$ Residual\\[0.5em]
    $\hat{\bm{y}}_z = \hat{\bm{y}}^{\text{base}}_z + \lambda_{\text{base}} \cdot \text{residual}$
    \vspace{1.5em}}}
    \caption{
    Overview of \method.
    A frozen base forecaster produces $\hat{\bm{y}}^{\text{base}}_z$ in z-space.
    The delta module, conditioned on both the time series history and LLM-refined news regime signals, predicts a structured residual decomposed into slow (level shift), shape (intra-day pattern), and spike (extreme event) components.
    Three modulation knobs---$\lambda_{\text{base}}$, $g_{\text{shape}}$, and $b_{\text{spike}}$---controlled by the regime representation, multiplicatively adjust these components.
    When no relevant news exists (relevance mass $m \leq \theta$), all knobs collapse to identity, and the prediction degrades gracefully to the base forecast.
    }
    \label{fig:architecture}
\end{figure}


\subsection{Stage 1: Base Backbone}
\label{sec:base}

The base model maps a normalized history window $\bm{h}_z \in \mathbb{R}^{L}$ to a prediction $\hat{\bm{y}}^{\text{base}}_z \in \mathbb{R}^{H}$ and an internal representation $\bm{f}_{\text{base}} \in \mathbb{R}^{D_{\text{base}}}$, where $L$ is the look-back length and $H$ is the forecast horizon.
We apply robust quantile normalization: $h_z = (h - \tilde{\mu}) / \tilde{\sigma}$, where $\tilde{\mu}$ is the training-set median and $\tilde{\sigma} = \text{IQR} / 1.349$, which provides robustness to price spikes.
The base is trained with Smooth-L1 loss in z-space (or a winsorized quantile loss for price) and is subsequently \emph{frozen} for all downstream usage.
The backbone is chosen per dataset from a pool of PatchTST~\citep{PLACEHOLDER_patchtst_nie2023}, DLinear~\citep{PLACEHOLDER_dlinear_zeng2023}, NLinear, and MLP variants (cf.\ Section~\ref{sec:setup}).


\subsection{News Distillation Pipeline}
\label{sec:news}

Raw news articles are transformed into a structured representation through two offline stages.

\paragraph{Schema refinement.}
Each article is processed by an LLM with a structured prompt that extracts:
(i)~an actionability flag (non-actionable articles are permanently discarded);
(ii)~a 5-dimensional regime vector $\bm{r} \in [-1,1]^5$ covering market tightness, demand outlook, renewable surplus, volatility tone, and policy status;
(iii)~multi-label topic tags from a closed vocabulary of 15 categories (\eg \texttt{heatwave}, \texttt{supply\_tight}, \texttt{outage});
(iv)~an estimated influence horizon in days;
(v)~a self-assessed confidence score $c \in [0,1]$.
This schema is domain-specific but fixed once designed, requiring no per-article human annotation.

\paragraph{Regime bank construction.}
The refined articles are aggregated into a daily regime bank $\mathcal{B}$.
For each calendar day $d$, all articles with $\text{published\_at} \leq d \leq \text{published\_at} + \text{horizon\_days}$ and $\text{is\_actionable} = \text{true}$ are collected.
Each article receives a time-decayed weight $w = c \cdot \exp(-\text{age} / \tau)$, where $\tau = 5$ days.
The weighted average produces a daily regime vector $\bar{\bm{r}}_d$, topic mass $\bar{\bm{t}}_d \in \mathbb{R}^{15}$, and text embedding $\bar{\bm{e}}_d \in \mathbb{R}^{384}$ (from E5-small-v2~\citep{PLACEHOLDER_e5_wang2024}).
Temporal smoothing via exponential moving average ($\alpha = 0.5$) produces the final bank entries.
A \emph{relevance mass} $m_d = \sum_i w_i$ (unnormalized) quantifies the aggregate signal strength for day $d$; days with $m_d \leq \theta$ are treated as news-inactive.

\subsection{Stage 2: Delta Module}
\label{sec:delta}

The delta module predicts a structured residual that corrects the base forecast.
It consists of four components: a time-series encoder, a regime projector, three residual heads, and three modulation heads.

\paragraph{Time-series encoder.}
We employ a PatchTST-style encoder~\citep{PLACEHOLDER_patchtst_nie2023}.
The history $\bm{h}_z$ is split into overlapping patches of length $P{=}8$ with stride $S{=}4$, yielding $N_p = 1 + \lfloor(L-P)/S\rfloor$ patches.
Each patch is linearly projected to $D{=}128$ dimensions and augmented with learnable calendar embeddings (day-of-week, half-hour-of-day, holiday).
The frozen base representation $\bm{f}_{\text{base}}$ is prepended as an additional token.
A 2-layer Transformer encoder with 4 heads processes the sequence, producing per-token features $\bm{T} \in \mathbb{R}^{(N_p+1) \times D}$ and a mean-pooled summary $\bm{s} \in \mathbb{R}^{D}$.

\paragraph{Regime projector.}
The concatenation of $\bar{\bm{r}}$, $\bar{\bm{t}}$, and $\bar{\bm{e}}$ (total 404 dimensions) is projected through a two-layer MLP with GELU activations and LayerNorm to produce $\bm{v}_{\text{regime}} \in \mathbb{R}^{D}$.
An active mask $a = \mathbf{1}[m > \theta]$ gates this representation: when news is inactive, $\bm{v}_{\text{regime}} = \bm{0}$.

\paragraph{Three-head residual decomposition.}
The time-series features $\bm{T}$ and $\bm{s}$ are processed by three specialized heads:
\begin{align}
    \text{Slow head:} \quad & s_{\text{slow}} = f_{\text{slow}}(\bm{s}) \in \mathbb{R}, \label{eq:slow} \\
    \text{Shape head:} \quad & \bm{q}_{\text{shape}} = f_{\text{shape}}(\bm{T}) - \overline{f_{\text{shape}}(\bm{T})} \in \mathbb{R}^{H}, \label{eq:shape} \\
    \text{Spike head:} \quad & \bm{q}_{\text{spike}}, \bm{g}_{\text{logit}} = f_{\text{spike}}(\bm{T}) \in \mathbb{R}^{H} \times \mathbb{R}^{H}, \label{eq:spike}
\end{align}
where the slow head captures level shifts (broadcast across the horizon), the shape head outputs a zero-mean intra-day pattern, and the spike head produces both spike amplitudes and gating logits.
Training targets for these heads are derived from the true residual $\bm{y}_z - \hat{\bm{y}}^{\text{base}}_z$ by decomposition against a pre-computed calendar baseline ($7 \times 48$ day-of-week $\times$ half-hour lookup table) and quantile-based spike thresholds.

\paragraph{Multiplicative modulation.}
The regime representation $\bm{v}_{\text{regime}}$ controls three bounded modulation knobs via the modulation heads:
\begin{align}
    \lambda_{\text{ts}} &= \lambda_{\min} + \lambda_{\text{ts}}^{\text{cap}} \cdot \sigma\!\left(f_{\text{trust}}(\bm{s})\right), \label{eq:lambda_ts} \\
    \Delta\lambda_{\text{news}} &= \lambda_{\text{news}}^{\text{cap}} \cdot \tanh\!\left(f_{\text{news}}(\bm{v}_{\text{regime}}) \cdot a\right), \label{eq:lambda_news} \\
    \lambda_{\text{base}} &= \text{clamp}\!\left(\lambda_{\text{ts}} + \Delta\lambda_{\text{news}},\; \lambda_{\min},\; \lambda_{\max}\right) \cdot \sigma_{\text{soft}}, \label{eq:lambda_base} \\
    g_{\text{shape}} &= 1 + g^{\text{cap}} \cdot \tanh\!\left(f_{\text{gain}}(\bm{v}_{\text{regime}}) \cdot a\right), \label{eq:shape_gain} \\
    b_{\text{spike}} &= b^{\text{cap}} \cdot \tanh\!\left(f_{\text{spike\_prior}}(\bm{v}_{\text{regime}}) \cdot a\right), \label{eq:spike_bias}
\end{align}
where $\sigma_{\text{soft}} = \sigma((m - \theta) / \tau_{\text{soft}})$ is a soft gate that smoothly shrinks $\lambda_{\text{base}}$ toward $\lambda_{\min}$ for low-relevance samples.
The final prediction is:
\begin{equation}
    \hat{\bm{y}}_z = \hat{\bm{y}}^{\text{base}}_z + \lambda_{\text{base}} \cdot \Big(\underbrace{s_{\text{slow}} \cdot \bm{1}}_{\text{level}} + \underbrace{g_{\text{shape}} \cdot \bm{q}_{\text{shape}}}_{\text{pattern}} + \underbrace{\sigma(\bm{g}_{\text{logit}} + b_{\text{spike}}) \odot \bm{q}_{\text{spike}}}_{\text{spikes}}\Big).
    \label{eq:final_pred}
\end{equation}

\paragraph{Degradation guarantee.}
When no news is active ($a{=}0$), we have $\Delta\lambda_{\text{news}}{=}0$, $g_{\text{shape}}{=}1$, and $b_{\text{spike}}{=}0$. The soft gate further pushes $\lambda_{\text{base}} \to \lambda_{\min}$.
Thus the worst-case deviation from the base forecast is bounded:
\begin{equation}
    \| \hat{\bm{y}}_z - \hat{\bm{y}}^{\text{base}}_z \|_\infty \leq \lambda_{\min} \cdot \max_h |r_h|,
    \label{eq:bound}
\end{equation}
where $r_h$ is the $h$-th element of the TS-only residual.
With $\lambda_{\min}{=}0.05$, the delta contribution is at most 5\% of the residual magnitude---a formal safety property that additive fusion methods cannot provide.


\subsection{Training Procedure}
\label{sec:training}

\paragraph{Two-phase delta training.}
Before joint training, a \emph{pretrain phase} (12~epochs) trains auxiliary regression heads on the regime representation to predict volatility ratio, spike flag, and shape deviation from the base residual.
The learned first-layer weights are transferred to initialize the modulation heads, providing a warm start for the trust, shape gain, and spike bias MLPs.

\paragraph{Loss function.}
The total loss combines six terms:
\begin{equation}
    \mathcal{L} = \mathcal{L}_{\text{point}} + w_s \mathcal{L}_{\text{slow}} + w_q \mathcal{L}_{\text{shape}} + w_k (\mathcal{L}_{\text{spike}} + w_g \mathcal{L}_{\text{gate}}) + w_c \mathcal{L}_{\text{consist}} + w_f \mathcal{L}_{\text{counter}},
    \label{eq:loss}
\end{equation}
where $\mathcal{L}_{\text{point}}$ is Smooth-L1 in z-space (load/gas) or winsorized raw-space with quantile loss (price), $\mathcal{L}_{\text{slow/shape}}$ are MSE on the decomposed targets, $\mathcal{L}_{\text{spike}}$ is BCE on spike gate activations, $\mathcal{L}_{\text{consist}}$ penalizes deviation between news-on and news-off predictions on news-\emph{inactive} samples (enforcing the degradation property during training), and $\mathcal{L}_{\text{counter}}$ encourages orthogonality between the TS-only and news-modulated residual components.
During training, 30\% of batches have their regime pack randomly blanked (set to zero), acting as a news-dropout regularizer.

\paragraph{Hard-sample mining.}
An importance sampler upweights batches where the base model has large residuals (top 10\% by magnitude, with 60\% probability of being selected), focusing the delta module's capacity on samples where correction is most needed.


%% ============================================================
%% 4. COUNTERFACTUAL EVALUATION PROTOCOL
%% ============================================================
\section{Counterfactual Evaluation Protocol}
\label{sec:counterfactual}

A central concern with text-augmented forecasting is whether improvements stem from genuine causal news signals or from statistical overfitting to the news-time-series correlation structure.
We address this with a three-condition evaluation protocol applied to every test sample:

\begin{enumerate}[leftmargin=1.5em, itemsep=2pt]
    \item \textbf{Normal}: Full regime pack from the regime bank.
    \item \textbf{News-blanked}: Regime pack set to all zeros, simulating complete absence of news. For \emph{news-inactive} samples ($m \leq \theta$), this should yield nearly identical predictions to the normal condition; any gap indicates unwanted news leakage.
    \item \textbf{Date-permuted}: Regime pack entries are shuffled across dates (with fixed seed for reproducibility), breaking the temporal alignment between news and time series. For \emph{news-active} samples, this should degrade performance if the model is genuinely using news content rather than just the presence/absence of news.
\end{enumerate}

We report four diagnostic metrics:
(i)~\textbf{inactive blank gap}: $|(\text{MAE}_{\text{normal}} - \text{MAE}_{\text{blank}}) / \text{MAE}_{\text{blank}}|$ on inactive samples (target: $< 2\%$);
(ii)~\textbf{active blank gain}: $\text{MAE}_{\text{blank}} - \text{MAE}_{\text{normal}}$ on active samples in z-space (positive $=$ news helps);
(iii)~\textbf{active permuted gain}: $\text{MAE}_{\text{permuted}} - \text{MAE}_{\text{normal}}$ on active samples (positive $=$ temporal alignment matters);
(iv)~\textbf{modulation diagnostics}: distributions of $\lambda_{\text{base}}$, $g_{\text{shape}}$, and $b_{\text{spike}}$ stratified by active/inactive status.


%% ============================================================
%% 5. EXPERIMENTS
%% ============================================================
\section{Experiments}
\label{sec:experiments}

\subsection{Datasets}
\label{sec:datasets}

We evaluate on five datasets spanning different domains, frequencies, and distributional characteristics (Table~\ref{tab:datasets}).
Forecast horizons are reported in \emph{steps} of the native sampling frequency (e.g., $H{=}48$ corresponds to one day for 30-min load/price and two days for hourly gas/traffic/Bitcoin).

\begin{table}[t]
    \centering
    \caption{Dataset characteristics. Tail severity ranges from moderate (near-Gaussian) to heavy (extreme right skew with $\sigma \approx 635$). Horizons list the set evaluated in this paper.}
    \label{tab:datasets}
    \small
    \begin{tabular}{lccccll}
        \toprule
        \textbf{Dataset} & \textbf{Freq.} & \textbf{$H$} & \textbf{Unit} & \textbf{Tail} & \textbf{Base} & \textbf{News Source} \\
        \midrule
        NSW Load & 30 min & 48--720 & MW & Moderate & PatchTST & Curated energy (Free News, RenewEconomy, PV-Mag, WattClarity) \\
        NSW Traffic & 1 hour & 48--720 & count & Moderate & PatchTST & GDELT 2024 (global) \\
        Gas Demand (NL) & 1 hour & 48 & -- & Moderate & DLinear & CBS + Free News archival \\
        Bitcoin & 1 hour & 48 & USD & Moderate & MLP & GDELT 2024 (global) \\
        NSW Price & 30 min & 48 & \$/MWh & Heavy & PatchTST/NLinear & Curated energy (as NSW Load) \\
        \bottomrule
    \end{tabular}
\end{table}

\textbf{NSW Load and Price.} Australian National Electricity Market data from the New South Wales region, sourced from AEMO for 2024.
The load series exhibits strong calendar regularity (weekday/weekend, morning/evening peaks); the price series has extreme right-skewed distribution with spikes reaching \$5{,}000--15{,}000/MWh against a median of ${\sim}$\$50--100/MWh.
Both datasets share a curated news corpus scraped from Free News, RenewEconomy, PV-Magazine, and WattClarity, with coverage of NEM policy, outages, and weather-driven demand.

\textbf{NSW Traffic.} Hourly traffic counts for 2024 from NSW permanent count stations, aggregated into a single hourly series.
Strong daily and weekly seasonality with suppressed-volume periods around holidays and weather events.
We pair this series with the open GDELT 2024 global news corpus, filtered to English-language items, to test whether a \emph{non-curated} news source can still supply a useful regime signal.

\textbf{European Gas Demand.} Hourly natural gas consumption from the Netherlands in 2025, sourced from ENTSOG.
Strong seasonal and day-of-week patterns driven by heating demand.

\textbf{Bitcoin Hourly.} Hourly Bitcoin opening prices from 2024, paired with the GDELT 2024 corpus as a broad, cross-topic news stream rather than a hand-picked crypto feed.

\subsection{Experimental Setup}
\label{sec:setup}

\paragraph{Base backbone.}
We select the base backbone per dataset from a pool of PatchTST, DLinear, NLinear, and MLP variants: PatchTST for NSW load and NSW traffic; DLinear for gas demand; MLP for Bitcoin; and PatchTST/NLinear for NSW price.
Load, gas demand, Bitcoin, and price bases are trained for 40 epochs, while NSW traffic uses 100 epochs (patience 5); all bases use Smooth-L1 loss at lr=$10^{-3}$ (or winsorized quantile loss for price) and are \emph{frozen} for all downstream usage.

\paragraph{Delta module.}
All datasets share a common delta architecture: a PatchTST encoder with 2 Transformer layers, 4 attention heads, hidden size 128, patch length 8, and stride 4.
Text embeddings are produced by E5-small-v2 (384-d)~\citep{PLACEHOLDER_e5_wang2024}.
Training uses AdamW with gradient clipping at 1.0 and a 10\%-warmup cosine schedule; the pretrain phase runs for 12 epochs at lr=$10^{-3}$, followed by the main delta phase (30--50 epochs at lr=$10^{-4}$, except NSW price which uses $5{\times}10^{-6}$ and gas demand which uses $5{\times}10^{-6}$).
The default modulation caps are $\lambda_{\min}{=}0.05$, $\lambda_{\text{ts}}^{\text{cap}}{=}0.30$, $\lambda_{\text{news}}^{\text{cap}}{=}0.12$, $\lambda_{\max}{=}0.45$, $g^{\text{cap}}{=}0.20$, $b^{\text{cap}}{=}0.75$, with the active-mass threshold $\theta$ tuned per dataset (0.1 for gas demand, 0.2 for Bitcoin, 0.7 for load/traffic/price).
For Bitcoin and long-horizon runs, the caps are further annealed with the forecast horizon (see Appendix~\ref{app:hyperparams}).

\paragraph{Evaluation metrics.}
(i)~MSE and MAE in original units;
(ii)~\textbf{Skill score}: $1 - \text{MAE}_{\text{delta}} / \text{MAE}_{\text{base}}$ (positive $=$ improvement);
(iii)~\textbf{Help rate}: fraction of test samples where $|\text{err}_{\text{delta}}| < |\text{err}_{\text{base}}|$;
(iv)~\textbf{Top-10\% help rate}: help rate restricted to the 10\% of samples with largest base residuals;
(v)~Counterfactual diagnostics from Section~\ref{sec:counterfactual}.


\subsection{Main Results}
\label{sec:main_results}

Table~\ref{tab:main_results} presents the main comparison across all datasets.
\emph{[Placeholder: numerical entries will be populated once the canonical sweep over all datasets and horizons has finished under the current framework configuration.]}

\begin{table}[t]
    \centering
    \caption{Main results. Skill Score = $1 - \text{MAE}_\text{delta} / \text{MAE}_\text{base}$. Help Rate and Top-10\% Help Rate measure the fraction of samples improved overall and in the hardest decile. $\bar{\lambda}$ is the mean $\lambda_{\text{base}}$ across test samples. Regime\% is the fraction of samples with active news. Best skill score per dataset in \textbf{bold}.}
    \label{tab:main_results}
    \small
    \begin{tabular}{llcccccccc}
        \toprule
        \textbf{Dataset} & \textbf{Base} & \textbf{$H$} & \textbf{Base MAE} & \textbf{Delta MAE} & \textbf{Skill$\uparrow$} & \textbf{Help$\uparrow$} & \textbf{Top-10\%$\uparrow$} & $\bar{\lambda}$ & Regime\% \\
        \midrule
        NSW Load    & PatchTST & 48  & -- & -- & -- & -- & -- & -- & -- \\
        NSW Load    & PatchTST & 192 & -- & -- & -- & -- & -- & -- & -- \\
        NSW Load    & PatchTST & 720 & -- & -- & -- & -- & -- & -- & -- \\
        \midrule
        NSW Traffic & PatchTST & 48  & -- & -- & -- & -- & -- & -- & -- \\
        NSW Traffic & PatchTST & 192 & -- & -- & -- & -- & -- & -- & -- \\
        NSW Traffic & PatchTST & 720 & -- & -- & -- & -- & -- & -- & -- \\
        \midrule
        Gas Demand  & DLinear  & 48  & -- & -- & -- & -- & -- & -- & -- \\
        \midrule
        Bitcoin     & MLP      & 48  & -- & -- & -- & -- & -- & -- & -- \\
        \midrule
        NSW Price   & PatchTST & 48  & -- & -- & -- & -- & -- & -- & -- \\
        \bottomrule
    \end{tabular}
\end{table}

\paragraph{NSW Load.}
With the PatchTST base and a curated energy-news corpus, the load series is the prototypical ``short causal chain'' case for \method: heatwave warnings and outage notices translate almost directly into next-day demand shifts.
We expect positive skill scores at short and medium horizons and strong top-decile help rates, tapering off at $H{=}720$ where the 30-day-ahead signal in the news bank is inherently weaker.

\paragraph{NSW Traffic.}
The traffic series is our \emph{non-curated news} stress test: the news side of the pipeline uses the raw GDELT 2024 corpus rather than a hand-picked feed, so a significant fraction of retrieved articles is only tangentially related to Sydney-area traffic.
This configuration tests whether the regime bank's LLM-refinement and confidence-weighted aggregation can extract a usable signal from an open, high-volume news stream.

\paragraph{Gas demand.}
Netherlands gas demand has a strong weather-driven seasonal component, and the CBS archival corpus supplies policy / storage / cold-snap context at daily granularity.
The DLinear base captures trend and seasonality cheaply, leaving \method's delta module to modulate shape and spike components under cold-snap regimes.

\paragraph{Bitcoin.}
Bitcoin's hourly price dynamics combine weak calendar structure with sparse but occasionally high-impact news.
We expect the low-$\lambda$, low-regime\% regime of the framework---a mostly-TS correction with occasional news nudges---to dominate.

\paragraph{NSW Price (instructive failure).}
Even when the delta module improves a majority of samples (help rate $>50\%$), overall MAE can worsen on the price series.
Section~\ref{sec:analysis} formalizes this failure mode as \emph{help-rate/MAE decoupling}.


\subsection{Counterfactual Analysis}
\label{sec:counterfactual_results}

Table~\ref{tab:counterfactual} reports the counterfactual diagnostics across all datasets.
\emph{[Placeholder: numerical entries pending the final pass of the counterfactual protocol on the current configurations.]}

\begin{table}[t]
    \centering
    \caption{Counterfactual evaluation across datasets. Inactive Blank Gap measures prediction consistency when news is blanked on news-inactive samples (target: $<2\%$). Active Blank$-$Normal and Active Permuted$-$Normal report MAE changes in z-space on news-active samples; positive values indicate that correct, time-aligned news helps. Modulation diagnostics check that $\bar{\lambda}_{\text{base}}$ and spike gate activations fall in acceptable ranges.}
    \label{tab:counterfactual}
    \small
    \begin{tabular}{lccccc}
        \toprule
        \textbf{Dataset} & \textbf{Inactive Blank Gap} & \textbf{Active Blank$-$Normal} & \textbf{Active Perm$-$Normal} & $\bar{\lambda}_{\text{base}}$ & \textbf{Spike Gate Hit Rate} \\
        \midrule
        NSW Load    & -- & -- & -- & -- & -- \\
        NSW Traffic & -- & -- & -- & -- & -- \\
        Gas Demand  & -- & -- & -- & -- & -- \\
        Bitcoin     & -- & -- & -- & -- & -- \\
        NSW Price   & -- & -- & -- & -- & -- \\
        \bottomrule
    \end{tabular}
\end{table}

The protocol is designed to confirm three properties of a well-behaved \method deployment:
(1)~\textbf{No news leakage}: the inactive blank gap should be well within the 2\% tolerance, confirming that multiplicative gating silences news influence on news-inactive samples.
(2)~\textbf{News helps active samples}: using correct news should outperform news-blanked predictions on news-active samples.
(3)~\textbf{Temporal alignment matters}: shuffling news-date assignments should degrade predictions, confirming that improvements depend on news \emph{content} rather than merely its presence.


\subsection{Analysis: Help-Rate/MAE Decoupling}
\label{sec:analysis}

The NSW Price results reveal that a delta module can improve the \emph{majority} of samples (help rate $>50\%$) while \emph{worsening} overall MAE.
This occurs because MAE weights errors by magnitude while help rate counts samples.

Let $\epsilon^+_i = |e^{\text{base}}_i| - |e^{\text{delta}}_i|$ denote the error reduction for sample $i$.
Then:
\begin{equation}
    \text{MAE}_{\text{base}} - \text{MAE}_{\text{delta}} = \frac{1}{N}\sum_{i=1}^N \epsilon^+_i = \underbrace{\frac{|\mathcal{H}|}{N}\,\mathbb{E}[\epsilon^+ \mid \mathcal{H}]}_{\text{helped benefit}} - \underbrace{\frac{|\mathcal{U}|}{N}\,\mathbb{E}[|\epsilon^+| \mid \mathcal{U}]}_{\text{hurt cost}},
    \label{eq:decoupling}
\end{equation}
where $\mathcal{H}$ and $\mathcal{U}$ are the helped and hurt sample sets.

For price data ($\sigma \approx 635$), the distribution is extremely heavy-tailed: helped samples each gain \$2--5, but hurt samples---concentrated in spike periods---each lose \$50--200.
The asymmetry in conditional expectations causes the hurt cost to dominate despite $|\mathcal{H}| > |\mathcal{U}|$.
For load data ($\sigma$ much smaller), the conditional expectations are roughly symmetric, so help rate directly translates to MAE improvement.

We identify four structural reasons for this failure:
(i)~the three-head decomposition (slow/shape/spike) matches the regular calendar structure of load but not the erratic structure of price residuals;
(ii)~z-space normalization amplifies small prediction errors into large raw-space errors when $\sigma$ is large;
(iii)~the news$\to$price causal chain (news $\to$ load $\to$ supply/demand $\to$ price) is longer and noisier than news$\to$load;
(iv)~spike gate false positives on price have disproportionately large costs.


%% ============================================================
%% 6. ABLATION STUDIES
%% ============================================================
\section{Ablation Studies}
\label{sec:ablation}

We ablate along four axes: (i)~base backbone choice, (ii)~active-mass threshold $\theta$, (iii)~news corpus (curated vs.\ open GDELT), and (iv)~forecast horizon.
Result tables for each study are deferred to Appendix~\ref{app:ablations} and will be populated once the current sweep has completed.

\paragraph{Effect of base backbone.}
We compare PatchTST, DLinear, NLinear, and MLP as the frozen base on NSW load, NSW traffic, and NSW price.
The interest is whether a stronger base starves the delta of useful residual signal (a risk flagged in our internal analysis) or whether the multiplicative modulation remains additive even on top of a Transformer base.

\paragraph{Effect of active-mass threshold $\theta$.}
The regime bank gates news influence by relevance mass $m_d$; moving $\theta$ changes the fraction of samples considered news-active.
Lower $\theta$ exposes the delta to more (noisier) news, while higher $\theta$ keeps only high-confidence regime days.
We sweep $\theta \in \{0.1, 0.2, 0.3, 0.7\}$ across datasets.

\paragraph{Effect of news corpus.}
On NSW load and NSW traffic we compare a \emph{curated} dataset-specific corpus (Free News / RenewEconomy / PV-Magazine / WattClarity for load) against the \emph{open} GDELT 2024 global corpus, holding all other hyperparameters fixed.
This ablation isolates the value of domain-specific news curation vs.\ the robustness of the LLM-refinement pipeline to noisy sources.

\paragraph{Effect of forecast horizon.}
For NSW load and NSW traffic we sweep $H \in \{48, 96, 192, 336, 720\}$ and anneal the modulation caps $\{\lambda_{\max}, \lambda_{\text{ts}}^{\text{cap}}, \lambda_{\text{news}}^{\text{cap}}, g^{\text{cap}}, b^{\text{cap}}\}$ with the horizon (see Appendix~\ref{app:hyperparams}).
The expected pattern is that skill score declines with horizon because the news regime bank's effective causal range is days rather than weeks.


%% ============================================================
%% 7. LIMITATIONS
%% ============================================================
\section{Limitations}
\label{sec:limitations}

\textbf{Heavy-tailed distributions.}
As demonstrated by the NSW Price results, \method struggles with extremely heavy-tailed series where rare spike errors disproportionately impact MAE.
The bounded modulation knobs---while providing safety---limit the delta module's ability to make the large corrections needed during price spikes.

\textbf{Causal chain length.}
The framework is most effective when news has a short causal chain to the target variable (weather $\to$ load, weather $\to$ gas demand).
For variables at the end of longer chains (news $\to$ load $\to$ supply/demand $\to$ price), the regime signal degrades.

\textbf{Daily granularity.}
The regime bank operates at daily granularity, potentially missing intra-day events critical for high-frequency price forecasting.

\textbf{LLM dependency.}
The offline refinement step requires LLM API access.
While this is a one-time cost per corpus and does not affect inference latency, it introduces a dependency.

\textbf{Single-series evaluation.}
Our experiments focus on univariate forecasting.
Extension to multivariate settings remains future work.


%% ============================================================
%% 8. CONCLUSION
%% ============================================================
\section{Conclusion}
\label{sec:conclusion}

We presented \method, a two-stage framework for incorporating LLM-refined news signals into time series forecasting through multiplicative residual modulation.
The key design choice---news signals can only \emph{scale} time-series-derived residual components, never inject additive predictions---provides a formal degradation bound and prevents the model from hallucinating forecasts based on text alone.
Our counterfactual evaluation protocol with news-blanking and date-permutation tests provides causal evidence that improvements are attributable to news content.
Experiments cover five datasets (NSW load, NSW traffic, Netherlands gas demand, Bitcoin, and NSW price), five forecast horizons up to 720 steps, and a mix of curated and open (GDELT) news corpora.
The heavy-tailed price series reveals a \emph{help-rate/MAE decoupling} failure mode that motivates future work on tail-aware modulation strategies and on base architectures that can selectively hand off residual structure to the delta module.

%% ============================================================
%% REFERENCES
%% ============================================================
\bibliographystyle{plainnat}
\bibliography{references}

%% ============================================================
%% APPENDIX
%% ============================================================
\newpage
\appendix

\section{Hyperparameter Details}
\label{app:hyperparams}

Table~\ref{tab:hyperparams} lists the complete hyperparameter settings.
Values reflect the current framework configuration; horizon-specific overrides for Bitcoin and long-horizon sweeps are applied via the run scripts.

\begin{table}[h]
    \centering
    \caption{Complete hyperparameter settings across datasets. Backbone choice is per dataset; all delta modules share the same PatchTST encoder.}
    \label{tab:hyperparams}
    \small
    \begin{tabular}{lccccc}
        \toprule
        \textbf{Parameter} & \textbf{Load} & \textbf{Traffic} & \textbf{Gas} & \textbf{Bitcoin} & \textbf{Price} \\
        \midrule
        Base backbone & PatchTST & PatchTST & DLinear & MLP & PatchTST/NLinear \\
        Base epochs / lr & 40 / $10^{-3}$ & 100 / $10^{-3}$ & 40 / $10^{-3}$ & 40 / $10^{-3}$ & 40 / $10^{-3}$ \\
        Early-stop patience & 5 & 5 & 5 & 5 & 5 \\
        \midrule
        Delta hidden / layers / heads & 128/2/4 & 128/2/4 & 128/2/4 & 128/2/4 & 128/2/4 \\
        Patch len / stride & 8/4 & 8/4 & 8/4 & 8/4 & 8/4 \\
        Delta lr / epochs & $10^{-4}$/30 & $10^{-4}$/50 & $5{\times}10^{-6}$/30 & $10^{-4}$/30 & $5{\times}10^{-6}$/30 \\
        Pretrain epochs / lr & 12/$10^{-3}$ & 12/$10^{-3}$ & 12/$10^{-3}$ & 12/$10^{-3}$ & 12/$10^{-3}$ \\
        Warmup / Grad clip & 10\%/1.0 & 10\%/1.0 & 10\%/1.0 & 10\%/1.0 & 10\%/1.0 \\
        \midrule
        $\lambda_{\min}$/$\lambda_{\text{ts}}^{\text{cap}}$/$\lambda_{\text{news}}^{\text{cap}}$/$\lambda_{\max}$ & .05/.30/.12/.45 & .05/.30/.12/.60 & .05/.30/.12/.45 & .05/.30/.12/.45 & .05/.30/.12/.45 \\
        $g^{\text{cap}}$ / $b^{\text{cap}}$ & 0.20/0.75 & 0.70/0.75 & 0.20/0.75 & 0.20--0.30/0.75 & 0.20--0.30/0.30--0.75 \\
        $\theta$ (active mass) & 0.7 & 0.7 & 0.1 & 0.2 & 0 (always active) \\
        Spike $k$ / target \% & 3.0/10\% & 3.0/10\% & 3.0/10\% & 3.0/10\% & 4.0/6\% \\
        News blank prob & 30\% & 30\% & 30\% & 30\% & 30\% \\
        Text encoder & E5-sm-v2 & E5-sm-v2 & E5-sm-v2 & E5-sm-v2 & E5-sm-v2 \\
        News corpus & Curated energy & GDELT 2024 & CBS + Free News & GDELT 2024 & Curated energy \\
        \bottomrule
    \end{tabular}
\end{table}

\section{Per-Dataset Ablation Tables}
\label{app:ablations}

% [Placeholder: per-dataset tables for backbone, $\theta$, news-corpus, and horizon ablations
% will be populated once the corresponding sweeps are run under the current framework.]

\section{Additional Counterfactual Results}
\label{app:counterfactual}

% [Placeholder: full counterfactual diagnostics (per-dataset, per-horizon) will be added here.
% Headline summary is provided in Table~\ref{tab:counterfactual}.]

\section{NeurIPS Paper Checklist}
\label{app:checklist}

% TODO: Complete NeurIPS paper checklist before submission

\end{document}
